\begin{abstract}
Providers of neural text-to-speech need to substantiate provenance claims about
audio they generated, which motivates watermarking the intermediate mel spectrogram.
Existing mel-domain schemes encode each bit as an absolute perturbation pattern and
degrade sharply under strong additive noise. We propose \texttt{RAWMER}, which
instead encodes each bit as the sign of an energy difference between two keyed groups
of mel bins, repeated across time blocks. Because both groups lie in the same local
time--frequency region, distortions affecting them equally---gain changes, mild
filtering, vocoder bias---cancel in the difference; a selection step uses the stored
clean reference to choose bin groups that avoid the dynamic-range boundary, carry
stable energy and are balanced in reliability. Against our re-implementation of
MelShield at matched clean quality (PESQ\,$\approx$\,3.51) on HiFi-GAN,
\texttt{RAWMER} raises bit accuracy under additive noise at 10 and 5\,dB SNR from
71.3\% and 63.3\% to 94.1\% and 85.9\%, with consistent gains on DiffWave. Further
experiments show chance-level false positives, $\approx$\,22\,KB references and
verification from unaligned 25\% clips. Time--frequency desynchronization remains the
main limitation and the most likely adaptive attack.
\end{abstract}
%
\begin{keywords}
Speech watermarking, provenance verification, mel spectrogram, neural
text-to-speech, additive-noise robustness
\end{keywords}
